\documentclass[../main.tex]{subfiles}
\begin{document}
%%%%%%%%%%%%%%%%%%%%%%%%%%%%%%%%%%%%%%%%%%%%%%%%%%%%%%%%%%%%%%%%%%%%%%%%%%%%%%%%%%%%%%%%%%%%%%%%%%%
\chapter{Opening and Importing Files in OmniGraffle}

This chapter provides you with a quick overview of the file and image types supported by
OmniGraffle. To help make it easy for you to find the file type you're looking for, we've
listed them in a convenient table, along with notations about whether that type is supported
by OmniGraffle Standard or Pro.

\begingroup
\small
\setlength{\LTleft}{0.075\textwidth}
\setlength{\LTright}{0.075\textwidth}
\renewcommand{\arraystretch}{1.25}
\begin{xltabular}{0.85\textwidth}{|>{\raggedright\arraybackslash}X|c|c|>{\raggedright\arraybackslash}X|}
  \caption{List of supported file types}\\
  \hline
  \rowcolor{ogheaderblue}
  \textcolor{white}{\textbf{File Type}} & \textcolor{white}{\textbf{Standard}} &
  \textcolor{white}{\textbf{Pro}} & \textcolor{white}{\textbf{How to Open/Import}} \\
  \hline
  \endhead
  \hline
  \endfoot
  PNG (\textit{.png}) & \ogyes & \ogyes & Drag onto a shape, layer, or canvas, or into the
  image wells provided in the \textbf{Image} and \textbf{Background Image} inspectors. \\
  \hline
  JPEG (\textit{.jpg}, \textit{.jpeg}) & \ogyes & \ogyes &  \\
  \hline
  GIF (\textit{.gif}) & \ogyes & \ogyes &  \\
  \hline
  BMP (\textit{.bmp}) & \ogyes & \ogyes &  \\
  \hline
  TIFF (\textit{.tiff}) & \ogyes & \ogyes &  \\
  \hline
  EPS (\textit{.eps}) & \ogyes & \ogyes &  \\
  \hline
  SVG (\textit{.svg}) & \ogyes & \ogyes &  \\
  \hline
  PDF (\textit{.pdf}) & \ogyes & \ogyes & Use one of the methods above to add a \textit{.pdf}
  file to the canvas as an image. Or, to add its individual elements as OmniGraffle objects,
  choose \textbf{File} $\blacktriangleright$ \textbf{Open} or drag the \textit{.pdf} file
  onto OmniGraffle's dock icon. \\
  \hline
  OmniGraffle Document (\textit{.graffle}) & \ogyes & \ogyes & Choose \textbf{File}
  $\blacktriangleright$ \textbf{Open} \\
  \hline
  OmniGraffle Stencil (\textit{.gstencil}) & \ogyes & \ogyes & Choose \textbf{File}
  $\blacktriangleright$ \textbf{Open} \\
  \hline
  OmniGraffle Template (\textit{.gtemplate}) & \ogyes & \ogyes & Choose \textbf{File}
  $\blacktriangleright$ \textbf{Open} \\
  \hline
  macOS Folder Structures & \ogyes & \ogyes & Drag a folder from the Finder onto OmniGraffle's
  app icon in the Dock \\
  \hline
  Framework Files (\textit{.framework}) & \ogyes & \ogyes & Drag the \textit{.framework} file
  from the Finder onto OmniGraffle's app icon in the Dock \\
  \hline
  Graphiz DOT Files (\textit{.dot}) & \ogyes & \ogyes & Drag the \textit{.dot} file onto
  OmniGraffle's app icon in the Dock \\
  \hline
  Adobe Photoshop (\textit{.psd}) & \ogyes & \ogyes & From the Finder, select and drag the
  \textit{.psd} file onto the canvas \\
  \hline
  Xcode Files (\textit{.xcodeproj}) & \ogno & \ogyes & Drag the \textit{.xcodeproj} file from
  the Finder onto OmniGraffle's app icon in the Dock \\
  \hline
  OmniOutliner (\textit{.oo3}) & \ogno & \ogyes & Choose \textbf{File} $\blacktriangleright$
  \textbf{Open} \\
  \hline
  Microsoft Visio (\textit{.vsd}, \textit{.vsdx}, \textit{.vdx}, \textit{.vst}, \textit{.vtx},
  \textit{.vss}, \textit{.vsx}, \textit{.vssx}) & \ogno & \ogyes & Choose \textbf{File}
  $\blacktriangleright$ \textbf{Open}; select the Visio file, and then click \textbf{Open}
  \\
  \hline
\end{xltabular}
\endgroup

\section{Importing OmniOutliner Files}

OmniGraffle interprets the contents of an OmniOutliner file as a diagram. As such, you need to
choose a diagram style when opening an OmniOutliner file in OmniGraffle:

\begin{enumerate}
  \item After choosing the OmniOutliner file to import, the \textbf{Import Outline} dialog
    appears:

    \ogfigure{0.800}{res/og7_import01.png}
  \item Using the pop-up menu, choose a diagram style:

    \ogfigure{0.800}{res/og7_import02.png}
  \item When you click \textbf{OK}, OmniGraffle renders the contents of the outline as a
    diagram:

    \ogfigure{0.800}{res/og7_import03.png}
\end{enumerate}

\begin{ogpro}
  When importing OmniOutliner files in OmniGraffle Pro, use the table to assign objects in the
  diagram style to columns in an outline:

  \ogfigure{0.800}{res/og7_import04.png}

  If you want to show multiple columns of information for each item, use a diagram style that
  has grouped objects. You can also assign a column to the objects' notes, or choose whether
  to place the outline in a completely new document or in a new canvas on the open document.
\end{ogpro}

\subsection{Importing Microsoft Visio Files}

\begin{ogpro}
  When importing Microsoft\textsuperscript{\textregistered{}}
  Visio\textsuperscript{\textregistered{}} (\textit{.vsd}, \textit{.vsdx}, \textit{.vdx},
  \textit{.vst}, \textit{.vtx}, \textit{.vss}, \textit{.vsx} \textit{.vssx}), it's important to
  note that Microsoft\textsuperscript{\textregistered{}}
  Visio\textsuperscript{\textregistered{}} saves its files in an XML-based binary file format
  which can be read by OmniGraffle Pro.

  While OmniGraffle Pro can open the binary Visio file, you are prompted to save the file as an
  OmniGraffle Stencil file (\textit{.gstencil}). Once imported, you can edit and add to the
  stencil file as needed. If you need to use that stencil in Visio again, do the following:

  \begin{enumerate}
    \item Open the stencil file in OmniGraffle.
    \item Choose \textbf{File} $\blacktriangleright$ \textbf{Export}
      (\textbf{Option-Command-E}).
    \item Choose \textbf{Visio} as the export format at the top of the Export panel.
    \item Click \textbf{Export}.
  \end{enumerate}

  OmniGraffle exports the file to Visio's \textit{.vdx} file format.

  Also note that graphics in the \textbf{Microsoft\textsuperscript{\textregistered{}} Enhanced
  Metafile (EMF)} format cannot be imported.
\end{ogpro}

\subsection{Supported DOT Syntax}

OmniGraffle understands a subset of the DOT language used by
\href{http://www.graphviz.org}{Graphviz}. This isn't a full implementation of DOT, but it can
be quite useful:

\begingroup
\small
\setlength{\LTleft}{0.075\textwidth}
\setlength{\LTright}{0.075\textwidth}
\renewcommand{\arraystretch}{1.25}
\begin{xltabular}{0.85\textwidth}{|l|>{\raggedright\arraybackslash}X|}
  \caption{Supported DOT syntax}\\
  \hline
  \rowcolor{ogheaderblue}
  \textcolor{white}{\textbf{Command}} & \textcolor{white}{\textbf{Entity Types}} \\
  \hline
  \endhead
  \hline
  \endfoot
  URL (not on graphs) & edge/line, node/shape \\
  \hline
  arrowhead & edge/line \\
  \hline
  arrowtail & edge/line \\
  \hline
  color & edge/line, node/shape \\
  \hline
  dir & edge/line \\
  \hline
  fillcolor & node/shape \\
  \hline
  fixedsize & node/shape \\
  \hline
  fontcolor & edge/line, node/shape \\
  \hline
  fontname & edge/line, node/shape \\
  \hline
  fontsize & edge/line, node/shape \\
  \hline
  headlabel & edge/line \\
  \hline
  label & edge/line, node/shape \\
  \hline
  labelfontcolor & edge/line \\
  \hline
  labelfontname & edge/line \\
  \hline
  labelfontsize & edge/line \\
  \hline
  orientation & node/shape \\
  \hline
  rank & node/shape \\
  \hline
  rankdir & graph/page \\
  \hline
  shape & node/shape \\
  \hline
  style & edge/line, node/shape \\
  \hline
  taillabel & edge/line \\
  \hline
  width & node/shape \\
  \hline
\end{xltabular}
\endgroup

\begin{ognote}
  All \href{http://www.graphviz.org/content/arrow-shapes}{Graphviz arrow shapes} are supported,
  but a few do not correspond exactly to OmniGraffle arrow types. The \texttt{inv},
  \texttt{invdot}, \texttt{invodot}, and \texttt{invempty} arrow types are replaced with the
  closest available OmniGraffle line ending.
\end{ognote}

\subsection{Using LinkBack to Embed Content from Another Application}

LinkBack is a way of embedding content between applications in such a way that you can always go
back and edit the content with the software used to create it originally.

\begin{ognote}
  For information about LinkBack, and a list of applications that support it, see the
  \href{http://linkback.org}{LinkBack project page} or the
  \href{http://github.com/omnigroup/LinkBack}{LinkBack repository on GitHub}.
\end{ognote}

To embed OmniGraffle content in another LinkBack application, just copy and paste the content
like you would normally. The same goes for embedding other applications' content in
OmniGraffle diagrams.

To embed OmniGraffle content inside a single object in OmniGraffle, select all of the objects on
your canvas (\textbf{Command-A}), choose \textbf{Edit} $\blacktriangleright$ \textbf{Copy As}
$\blacktriangleright$ \textbf{PDF} from the menu bar, and then Paste (\textbf{Command-V}). The
objects you copied become a single image; double-clicking the object opens a new OmniGraffle
window in which you can edit the objects to your heart's content.

To edit OmniGraffle content embedded in another application:

\begin{itemize}
  \item Double-click the OmniGraffle content in the other application's document; an OmniGraffle
    window appears containing the content.
  \item Edit the content in the OmniGraffle window, just as you normally would.
  \item In OmniGraffle, choose \textbf{File} $\blacktriangleright$ \textbf{Save}; the content in
    the other application updates automatically.
  \item To edit other LinkBack applications' content embedded in an OmniGraffle diagram, just
    double-click it. When you save in the other application, the content in the OmniGraffle
    diagram updates.
\end{itemize}
%%%%%%%%%%%%%%%%%%%%%%%%%%%%%%%%%%%%%%%%%%%%%%%%%%%%%%%%%%%%%%%%%%%%%%%%%%%%%%%%%%%%%%%%%%%%%%%%%%%
\end{document}
